\documentclass[11pt]{article}

% ── Packages ──────────────────────────────────────────────────────────────────
% NeurIPS submission style. Falls back to 1-inch margins if neurips_2024.sty
% is unavailable in the build environment (the .sty file ships in paper/).
\usepackage{neurips_2024}
\usepackage[utf8]{inputenc}
\usepackage[T1]{fontenc}
\usepackage{amsmath}
\usepackage{amssymb}
\usepackage{amsfonts}
\usepackage{amsthm}
\usepackage{algorithm}
\usepackage{algorithmic}
\usepackage{booktabs}
\usepackage{enumitem}
\usepackage{listings}
\usepackage{xcolor}
\usepackage{pifont}
\usepackage{hyperref}
\usepackage{url}
\usepackage{natbib}

% ── Theorem environments ─────────────────────────────────────────────────────
\newtheorem{theorem}{Theorem}
\newtheorem{lemma}[theorem]{Lemma}
\newtheorem{proposition}[theorem]{Proposition}
\newtheorem{corollary}[theorem]{Corollary}
\newtheorem{definition}{Definition}
\newtheorem{remark}{Remark}

% ── Listings configuration ───────────────────────────────────────────────────
\lstset{
  language=Python,
  basicstyle=\ttfamily\small,
  keywordstyle=\color{blue}\bfseries,
  commentstyle=\color{gray}\itshape,
  stringstyle=\color{red!70!black},
  numberstyle=\tiny\color{gray},
  numbers=left,
  numbersep=5pt,
  breaklines=true,
  showstringspaces=false,
  frame=single,
  rulecolor=\color{black!30},
  backgroundcolor=\color{gray!5},
  tabsize=4,
  captionpos=b
}

% ── Convenience macros ───────────────────────────────────────────────────────
\newcommand{\cmark}{\ding{51}}
\newcommand{\xmark}{\ding{55}}

% ══════════════════════════════════════════════════════════════════════════════
\title{TSAM: Page-Typed Information-Flow Control for\\
Multi-Tenant LLM Serving}

\author{%
  Anonymous Author(s)\\
  Affiliation\\
  \texttt{anon@example.com}%
}

\begin{document}
\maketitle

% ══════════════════════════════════════════════════════════════════════════════
% ABSTRACT
% ══════════════════════════════════════════════════════════════════════════════
\begin{abstract}
Multi-tenant serving of large language models batches requests from
distinct tenants into a single attention computation, creating a covert
data path from one tenant's KV-cache to another tenant's output.  Existing
mitigations either physically isolate tenants (5--10$\times$ memory cost,
no prefix sharing), project tenant keys into orthogonal subspaces (capped
at $\lfloor d_{\mathrm{model}}/d_{\mathrm{sub}} \rfloor$ tenants and
incompatible with shared prefixes), or apply ad-hoc attention masks
without a formal information-flow argument.  We present \textbf{TSAM}
(Tenant-Scoped Attention Masking), an information-flow-control mechanism
for PagedAttention~\citep{kwon2023vllm} that adds a one-byte
\textsc{Page\-Type} tag and a four-byte tenant identifier to each KV-cache
page and enforces a single hard-mask conditional inside the attention
kernel.  TSAM is the first design we are aware of that integrates
language-based information-flow control~\citep{denning1976lattice,
goguen1982security, sabelfeld2003language} with PagedAttention's page
abstraction, decoupling tenant scalability from model dimensionality
while preserving prefix sharing.

We give a typed model of KV-cache pages, prove
\emph{Attention-Channel Noninterference under Correct Page Typing}: under
four explicit deployment conditions (per-layer masking, clean shared
prefix, position-local non-attention sublayers, frozen shared weights),
the output at any tenant's positions is a deterministic function of only
that tenant's own-private and shared KV pairs across all $L$ transformer
layers.  Empirically, on a 24-layer transformer with random weights, TSAM
yields bitwise-identical attacker output regardless of the victim's
private tokens (max $\Delta = 0.0$ over 50 trials per depth), while a
linear probe extracts the victim's private prompt with $100\%$ accuracy
under vanilla causal attention versus random-chance ($11.7\%$ vs.\ $12.5\%$
baseline) under TSAM.  Per-page metadata adds $\approx$0.06\% of KV-cache
volume.  We are explicit about scope: TSAM addresses the
attention-channel only.  GPU memory side channels, malicious operators,
non-attention covert channels, and per-tenant LoRA adapters in the masked
path remain out of scope.  All code, the formal model, and the empirical
harness are open-source.
\end{abstract}

% ══════════════════════════════════════════════════════════════════════════════
% SECTION 1 — INTRODUCTION
% ══════════════════════════════════════════════════════════════════════════════
\section{Introduction}
\label{sec:introduction}

Large language models (LLMs) such as GPT-4~\citep{brown2020language,
achiam2023gpt4} power an expanding array of commercial applications---from
customer-facing chatbots to enterprise document analysis---that serve
thousands of tenants from a shared infrastructure.  The economic case is
compelling: a single A100 or H100 GPU can cost \$2--4/hour, and dedicating
one instance per tenant is infeasible beyond a handful of customers.
Systems such as vLLM~\citep{kwon2023vllm} address this through
\emph{PagedAttention}, which manages the KV-cache in fixed-size pages and
enables \emph{continuous batching} of requests from different tenants
within the same forward pass.  This design yields order-of-magnitude
improvements in GPU utilization and throughput.

However, batching requests from distinct tenants into a single attention
computation creates a subtle but critical vulnerability: tokens from
tenant~$t_i$ can, in principle, attend to KV-cache entries produced by
tenant~$t_j$'s private context.  Recent work has shown that LLMs can
memorize and regurgitate training data~\citep{carlini2021extracting}, and
analogous risks arise at inference time when private KV-cache pages are
accessible to cross-tenant queries.  The attention
mechanism~\citep{vaswani2017attention} computes
$\mathrm{softmax}(QK^\top/\sqrt{d_k})\,V$ over \emph{all} keys in the
batch; without explicit isolation, every query has a nonzero attention
weight on every key.

\paragraph{The gap.}
Despite the critical importance of multi-tenant LLM security, no existing
approach simultaneously achieves (1)~provably zero information leakage
between tenants, (2)~arbitrary tenant scalability, and (3)~compatibility
with prefix sharing---the technique that avoids redundant computation of
common system prompts.  This gap is a fundamental barrier to deploying
shared LLM serving in regulated industries (healthcare, finance, legal)
where formal isolation guarantees are a prerequisite.

\paragraph{Limitations of existing approaches.}
\emph{Process isolation} runs each tenant in a separate serving instance.
While trivially correct, this approach incurs 5--10$\times$ cost overhead
due to duplicated model weights and KV-cache, entirely negating the
economic advantage of shared serving.
\emph{Orthogonal subspace projection}~\citep{li2024orthogonal} maps each
tenant's keys into a mutually orthogonal subspace of dimension
$d_{\mathrm{sub}}$, ensuring that cross-tenant dot products vanish.
However, the number of orthogonal subspaces is bounded by
$d_{\mathrm{model}} / d_{\mathrm{sub}}$; for typical models with
$d_{\mathrm{model}} = 4096$ and $d_{\mathrm{sub}} = 256$, this yields a
hard ceiling of \textbf{16 tenants}.  Moreover, because each tenant's keys
are rotated into a private subspace, prefix sharing is fundamentally
destroyed: the same prompt produces different key vectors for different
tenants, and no two tenants can share a cache page.
\emph{Ad hoc attention masks} have been explored in concurrent
work~\citep{gao2024attentionguard}, but these approaches lack formal
information-flow proofs and do not address timing side channels.

\paragraph{Our approach.}
We present \textbf{TSAM} (Typed Serving with Tenant-Scoped Attention
Masking), a principled and lightweight extension to PagedAttention that
resolves all three limitations simultaneously.  TSAM introduces a
single-bit \emph{type tag} per KV-cache page---\texttt{shared} or
\texttt{private}($t$)---and enforces a hard-masking rule in the attention
kernel: query $q$ from tenant $t_i$ may attend to page $p$ if and only if
$\mathrm{type}(p) = \texttt{shared}$ or $\mathrm{owner}(p) = t_i$.  This
conditional adds fewer than 10 instructions per page in the inner loop of
FlashAttention-2 and yields a \emph{structural zero} in mutual
information:
\begin{equation}
  I\!\bigl(\mathrm{output}(q_i);\;\mathrm{KV}_{\mathrm{private}}(t_j)\bigr) = 0
  \quad \forall\; i \neq j.
  \label{eq:structural_zero}
\end{equation}
Because the mask is enforced at the page level rather than the subspace
level, tenant count is bounded only by available memory, and shared-prefix
pages remain tagged \texttt{shared}---preserving prefix sharing without
modification.

\paragraph{Contributions.}
We make the following contributions:
\begin{enumerate}[leftmargin=*,itemsep=2pt]
\item \textbf{Typed KV-cache pages.} We introduce a page-level type
  system for the KV-cache that cleanly separates shared and
  tenant-private state within the existing PagedAttention framework
  (\S\ref{sec:design}).

\item \textbf{Tenant-scoped attention masking.} We design a hard-masking
  rule enforced inside the attention kernel, yielding cross-tenant
  attention weights that are exactly zero by IEEE~754 arithmetic, with a
  single warp-uniform conditional per page (\S\ref{sec:design}).

\item \textbf{Attention-channel noninterference under correct page typing.}
  We prove that TSAM satisfies noninterference~\citep{goguen1982security}
  under the lattice model of information flow~\citep{denning1976lattice,
  sabelfeld2003language}, yielding $I = 0$ as a structural invariant
  rather than a statistical bound, and we extend the result to the full
  $L$-layer transformer stack under four explicit conditions (C1--C4)
  on per-layer masking, prefix cleanliness, position-locality, and
  weight-sharing (\S\ref{sec:formal}, \S\ref{sec:e2e}).

\item \textbf{Detection-triggered timing mitigation.} We design DAI, a
  per-tenant DFA that detects probing attacks via timing-variance and
  KL-divergence statistics and activates noise injection plus rate
  limiting once an attack is confirmed.  We are explicit that DAI is
  detection-triggered, not constant-time: pre-detection leakage is
  bounded but nonzero (\S\ref{sec:dai}).

\item \textbf{Reference implementation and empirical verification.} We
  release the full source: typed-page block manager, mask-generation
  code, Triton FlashAttention prototype, vLLM patch sketch, and the
  end-to-end noninterference test harness.  On a 24-layer transformer
  with random weights, TSAM yields bitwise-identical attacker output
  regardless of victim tokens, and a linear probing attack drops from
  $100\%$ (vanilla) to random chance $11.7\%$ (TSAM).
  Production GPU benchmarks (real vLLM fork, A100/H100 throughput,
  full Llama-2 inference) are explicitly out of scope of this version
  and are catalogued in \S\ref{sec:limitations} (\S\ref{sec:experiments}).
\end{enumerate}

\paragraph{Comparison with existing approaches.}
Table~\ref{tab:comparison} summarizes the landscape.  TSAM is the only
approach that achieves all four desiderata.

\begin{table}[t]
\centering
\caption{Comparison of multi-tenant LLM isolation approaches.}
\label{tab:comparison}
\vspace{4pt}
\small
\begin{tabular}{@{}lcccc@{}}
\toprule
\textbf{Approach} & \textbf{Tenant Capacity} & \textbf{Prefix Sharing} & \textbf{Attention-Channel Guarantee} & \textbf{Overhead} \\
\midrule
Process Isolation          & Memory-bound & \xmark & Trivial         & 5--10$\times$ \\
Orth.\ Projection~\citep{li2024orthogonal} & $\leq d_h / d_{\mathrm{sub}}$ & \xmark & $\varepsilon$-approx & Moderate \\
Ad hoc Masks~\citep{gao2024attentionguard}  & Memory-bound & \cmark & No formal proof  & Low \\
\textbf{TSAM (ours)}      & \textbf{Memory-bound} & \cmark & \textbf{$I{=}0$ structurally, under C1--C4} & \textbf{0.06\% metadata} \\
\bottomrule
\end{tabular}
\end{table}

% ══════════════════════════════════════════════════════════════════════════════
% SECTION 2 — BACKGROUND AND RELATED WORK
% ══════════════════════════════════════════════════════════════════════════════
\section{Background and Related Work}
\label{sec:background}

\subsection{Multi-Tenant LLM Serving}
\label{sec:bg_serving}

Modern LLM serving systems process requests from many tenants on shared
GPU infrastructure.  At the core of efficient serving lies the
\emph{key--value (KV) cache}: during autoregressive generation, the key
and value projections for all previously generated tokens are stored so
that attention need not be recomputed from scratch at each decoding step.
For a model with $L$ layers, $H$ attention heads, and head dimension
$d_h$, the KV-cache for a single sequence of length $n$ requires
$2 L H d_h n$ elements, consuming tens of gigabytes for long contexts.

\textbf{PagedAttention}~\citep{kwon2023vllm} addresses KV-cache memory
fragmentation by managing the cache in fixed-size \emph{pages} (typically
16 tokens each), analogous to virtual memory in operating systems.  Pages
are allocated on demand and can be shared across sequences via
copy-on-write, enabling \emph{prefix sharing}: when multiple requests
begin with the same system prompt, their common prefix pages are stored
once and referenced by all.

\textbf{Continuous batching}~\citep{yu2022orca} further improves
throughput by dynamically inserting and removing requests from the GPU
batch at each iteration, rather than waiting for an entire batch to
complete.  Together, PagedAttention and continuous batching allow a single
GPU to serve hundreds of concurrent sequences with high utilization.

\textbf{Prefix sharing} is especially important in multi-tenant settings
where many tenants share a common system prompt (e.g., ``You are a helpful
assistant'').  Without prefix sharing, the KV-cache for shared prompts
must be duplicated per tenant, wasting memory proportional to the number
of tenants times the prompt length.  Any isolation mechanism that destroys
prefix sharing therefore imposes a significant and scaling cost penalty.

\subsection{Information Flow Security}
\label{sec:bg_infoflow}

The formal study of information flow in computing systems originates with
Denning's \emph{lattice model}~\citep{denning1976lattice}, which assigns
security labels to data objects and defines permitted information flows via
a partial order on labels.  A system satisfies \emph{secure information
flow} if, for every pair of labels $\ell_1 \not\leq \ell_2$, no
information flows from objects at level $\ell_1$ to observers at level
$\ell_2$.

Goguen and Meseguer~\citep{goguen1982security} formalized this intuition
as \emph{noninterference}: a system is noninterfering with respect to a
high-security domain $H$ and a low-security domain $L$ if the outputs
visible to $L$ are identical regardless of the inputs from $H$.  Formally,
for a system function $f$ with inputs partitioned into $H$ and $L$
components:
\begin{equation}
  f(h_1, l) = f(h_2, l) \quad \forall\; h_1, h_2 \in H,\; \forall\; l \in L.
  \label{eq:noninterference}
\end{equation}
In the multi-tenant LLM setting, we treat each tenant's private data as a
separate high-security domain and require that the outputs for tenant
$t_i$ are independent of the private inputs of any other tenant $t_j$
($j \neq i$).

An equivalent information-theoretic characterization uses mutual
information~\citep{shannon1948mathematical, cover2006elements}: a system
achieves perfect isolation if
$I(\mathrm{output}_i;\;\mathrm{input}_j) = 0$ for all $i \neq j$.  When
this quantity is exactly zero (not merely bounded by $\epsilon$), we call
it a \emph{structural zero}---it holds by construction rather than by
parameter tuning or statistical estimation.

\subsection{Cache Side-Channel Attacks}
\label{sec:bg_sidechannel}

Hardware cache side-channel attacks exploit timing differences in memory
access patterns to infer secret data.
Osvik~et~al.~\citep{osvik2006cache} demonstrated \emph{access-driven}
attacks on AES, where the attacker observes which cache lines are evicted
during encryption.  Yarom and Falkner~\citep{yarom2014flush} introduced
\textsc{Flush+Reload}, a high-resolution attack exploiting shared memory
pages in last-level caches to achieve cross-core information leakage.

In the context of LLM serving, analogous side channels arise through the
\emph{KV-cache}.  When multiple tenants share a PagedAttention cache pool,
the set of pages allocated to a tenant, the timing of page allocations and
evictions, and the sequence lengths (which determine the number of pages
accessed during attention) can all leak information.  Specifically:
\begin{itemize}[leftmargin=*,itemsep=2pt]
  \item \textbf{Attention weight leakage.} Without masking, softmax
    attention assigns nonzero weight to every key in the batch, including
    keys from other tenants' private contexts.  Even small attention
    weights create a nonzero information channel.
  \item \textbf{Timing side channels.} The latency of attention
    computation depends on the total number of KV-cache pages accessed.
    An adversarial tenant can infer the sequence length of co-located
    tenants by measuring response latency, potentially revealing
    confidential information about the nature of their queries.
  \item \textbf{Cache collision attacks.} In shared GPU memory, two
    tenants accessing the same physical cache lines can observe
    interference patterns, analogous to classical cache collision attacks
    in CPUs.
\end{itemize}
TSAM addresses the first two channels directly: hard masking eliminates
attention weight leakage, and optional constant-time padding neutralizes
timing variation.  GPU-level cache collision attacks are orthogonal and
require hardware-level mitigations beyond the scope of this work.

\subsection{Existing Isolation Mechanisms}
\label{sec:bg_isolation}

We survey three categories of existing approaches to multi-tenant LLM
isolation.

\paragraph{Process isolation.}
The most straightforward approach runs a separate model instance per
tenant.  Each tenant's KV-cache is physically isolated, trivially
guaranteeing noninterference.  However, this approach requires duplicating
the full model weights ($\sim$14\,GB for a 7B-parameter model in FP16)
and KV-cache per tenant, yielding 5--10$\times$ memory overhead and
proportionally reduced throughput.  For $N$ tenants, the total memory
scales as $O(N \cdot (W + C))$ where $W$ is the model weight size and $C$
is the per-tenant cache size, compared to $O(W + N \cdot C)$ for shared
serving.  Process isolation is therefore economically viable only for a
small number of high-value tenants.

\paragraph{Orthogonal subspace projection.}
Li~et~al.~\citep{li2024orthogonal} propose projecting each tenant's key
vectors into mutually orthogonal subspaces of the attention head
dimension.  Given head dimension $d_h$, each tenant is assigned a
subspace of dimension $d_{\mathrm{sub}}$, and the key projection is
followed by a projection matrix $P_t \in \mathbb{R}^{d_h \times d_h}$
that zeros out all components outside tenant $t$'s subspace.  Since the
subspaces are orthogonal, cross-tenant dot products vanish:
\begin{equation}
  q_i^\top P_{t_i}^\top P_{t_j} k_j = q_i^\top (P_{t_i} P_{t_j})^\top k_j = 0
  \quad \text{for } i \neq j,
  \label{eq:orthogonal}
\end{equation}
because $P_{t_i} P_{t_j} = 0$ for orthogonal projections onto disjoint
subspaces.  This approach has two critical limitations:
\begin{enumerate}[leftmargin=*,itemsep=2pt]
  \item \textbf{Tenant scalability.} The maximum number of tenants is
    $d_h / d_{\mathrm{sub}}$.  For Llama-2's head dimension $d_h = 128$
    and a minimum practical subspace dimension $d_{\mathrm{sub}} = 8$,
    this yields at most \textbf{16 tenants}.  Reducing $d_{\mathrm{sub}}$
    further degrades model quality because the effective attention
    dimension per tenant shrinks.
  \item \textbf{Prefix sharing destruction.} Because each tenant's keys
    are multiplied by a different projection matrix $P_t$, the same input
    token produces different key vectors for different tenants.
    Consequently, no two tenants can share KV-cache pages for a common
    prefix---prefix sharing is \emph{fundamentally incompatible} with
    subspace projection.
\end{enumerate}

\paragraph{Ad hoc attention masks.}
Several concurrent works~\citep{gao2024attentionguard, rajput2024inference}
propose applying attention masks to prevent cross-tenant attention.  While
directionally similar to TSAM, these approaches differ in important ways:
they operate at the token level rather than the page level (incurring
higher metadata overhead), they do not provide formal information-flow
proofs or noninterference guarantees, and they do not address timing side
channels.  TSAM's contribution is to show that page-level typing, combined
with a formal security model grounded in the lattice framework of
Denning~\citep{denning1976lattice}, yields a provably correct and
practically efficient isolation mechanism.

% ==============================================================================
% SECTION 3 — THREAT MODEL
% ==============================================================================
\section{Threat Model}
\label{sec:threat}

We consider a multi-tenant LLM serving environment in which a shared inference engine (e.g., vLLM~\citep{kwon2023vllm}) processes requests from multiple tenants within a single batched forward pass.

\paragraph{Adversary model.}
The adversary $\mathcal{A}$ is a \emph{co-tenant}: a legitimate user of the serving platform who can submit arbitrary queries to the shared model. $\mathcal{A}$'s requests may be \emph{co-located} in the same continuous batch as a victim tenant $\mathcal{V}$'s requests, meaning that both tenants' KV-cache entries reside in GPU memory simultaneously.

We consider the \textbf{worst-case co-location scenario}: the serving engine employs prefix sharing, in which common prefixes are stored once and their KV-cache blocks are physically shared across tenants. In this setting, $\mathcal{A}$'s block table may reference $\mathcal{V}$'s physical blocks. We further assume $\mathcal{A}$ can measure end-to-end request latency with millisecond precision.

\paragraph{Adversary goals.}
\begin{enumerate}[leftmargin=*,itemsep=2pt]
    \item \textbf{Direct content extraction.} Extract information about $\mathcal{V}$'s private KV-cache content through attention-based information flow.
    \item \textbf{Presence detection.} Infer whether specific content exists in $\mathcal{V}$'s cache through timing side channels.
\end{enumerate}

\paragraph{Formal security goals.}
\begin{align}
    I\bigl(\mathrm{output}(q_\mathcal{A});\; \mathrm{KV}_{\mathrm{private}}(\mathcal{V})\bigr) &= 0, \label{eq:content-isolation} \\
    I\bigl(\mathrm{latency}(q_\mathcal{A});\; \mathrm{cache\_state}(\mathcal{V})\bigr) &\leq \varepsilon. \label{eq:timing-isolation}
\end{align}
Equation~\eqref{eq:content-isolation} demands \emph{perfect} content isolation enforced exactly by TSAM's attention masking. Equation~\eqref{eq:timing-isolation} bounds timing leakage to negligible $\varepsilon$, achieved through the DAI defense.

\paragraph{Out of scope.}
We exclude: (i) attacks targeting model weights or training data, (ii) physical side channels (power analysis, electromagnetic emanation), (iii) adversaries with direct GPU memory access, and (iv) covert channels through model output semantics.

% ==============================================================================
% SECTION 4 — TSAM SYSTEM DESIGN
% ==============================================================================
\section{TSAM: System Design}
\label{sec:design}

The fundamental insight behind TSAM is that \textbf{isolation is a data-flow property, not a representation property}. Prior work enforces isolation by transforming representations---projecting tenant keys into orthogonal subspaces---coupling isolation capacity to model dimensionality. TSAM instead enforces isolation by controlling data flow at the attention level, decoupling isolation from dimensionality entirely.

\subsection{Typed KV-Cache Pages}
\label{sec:design:pages}

We extend PagedAttention~\citep{kwon2023vllm} with a lightweight type system.

\begin{definition}[Page Type]
$\textsc{PageType} \in \{\textsc{Private}(0),\; \textsc{Shared}(1)\}$, encoded as a single-byte enumeration.
\end{definition}

\begin{definition}[Typed Page]
A typed page is a tuple $p = (\texttt{block\_id}, \texttt{tenant\_id}, \texttt{page\_type}, \texttt{access\_list})$ where:
\begin{itemize}[leftmargin=*,itemsep=2pt]
    \item \textbf{Private:} $\texttt{access\_list} = \{\texttt{tenant\_id}\}$. Only the owning tenant may attend.
    \item \textbf{Shared:} $\texttt{access\_list} = \mathcal{U}$ (universal). Any tenant may attend. Uses sentinel $\texttt{SHARED\_TENANT\_ID} = -1$.
\end{itemize}
\end{definition}

The \textsc{TypedPage} is a frozen dataclass with construction-time invariants: (i) \textsc{Private} pages require $\texttt{tenant\_id} \neq -1$ and $\texttt{access\_list} = \{\texttt{tenant\_id}\}$; (ii) \textsc{Shared} pages require $\texttt{tenant\_id} = -1$. Immutability prevents type-confusion attacks.

\paragraph{Metadata overhead.} Each page adds 5 bytes: 1 byte (\texttt{uint8} type) + 4 bytes (\texttt{int32} tenant ID). For 8,192 blocks, total metadata is $\approx$40\,KB ($\approx$0.06\% of KV-cache).

\subsection{TypedBlockManager}
\label{sec:design:manager}

The TypedBlockManager extends vLLM's block allocator with GPU-resident metadata tensors:
\begin{align}
    \texttt{page\_type} &: \texttt{uint8}[N_{\mathrm{blocks}}], \\
    \texttt{tenant\_id} &: \texttt{int32}[N_{\mathrm{blocks}}].
\end{align}

\paragraph{Allocation.} Two primitives: $\texttt{allocate\_private}(t, n)$ sets type$=0$, tenant$=t$; $\texttt{allocate\_shared}(n)$ sets type$=1$, tenant$=-1$. Shared blocks use reference counting. $\texttt{free\_tenant}(t)$ performs vectorized bulk cleanup.

\paragraph{Access mask.}
\begin{equation}
    \texttt{mask}[j] = \bigl(\texttt{page\_type}[j] = \textsc{Shared}\bigr) \;\lor\; \bigl(\texttt{tenant\_id}[j] = t_q\bigr).
\end{equation}
This is a single vectorized GPU operation, $O(N_{\mathrm{blocks}})$, independent of tenant count.

\subsection{TSAM Attention Masking}
\label{sec:design:attention}

For query token $q_i$ from tenant $t_q$ and key position $j$ in block $b_{\lfloor j/P \rfloor}$:
\begin{align}
    \mathrm{blocked}(j) &= \bigl(\tau(b_{\lfloor j/P \rfloor}) = \textsc{Private}\bigr) \;\land\; \bigl(t(b_{\lfloor j/P \rfloor}) \neq t_q\bigr), \\
    \mathrm{mask}(j) &= \lnot\, \mathrm{blocked}(j).
\end{align}
The mask is applied to attention scores before softmax:
\begin{equation}
    s'_{ij} = \begin{cases}
        q_i^\top k_j / \sqrt{d} & \text{if } \mathrm{mask}(j) = \textsc{True}, \\
        -\infty & \text{otherwise}.
    \end{cases}
\end{equation}
Under IEEE 754, $\exp(-\infty) = 0$ \emph{exactly}---a structural zero, not numerical approximation.

\begin{theorem}[TSAM Noninterference]
\label{thm:noninterference}
For distinct tenants $t_i, t_j$: $I\bigl(\mathrm{output}(q_i);\; \mathrm{KV}_{\textsc{Private}}(t_j)\bigr) = 0$.
\end{theorem}
\begin{proof}[Proof sketch]
Since $\alpha_{ij} = \exp(-\infty)/Z = 0$ for all blocked positions, the output $o_i = \sum_j \alpha_{ij} v_j$ has zero contribution from cross-tenant private values. The output is a deterministic function of only own-tenant and shared KV, hence $I = 0$ by the data processing inequality.
\end{proof}

\begin{algorithm}[t]
\caption{Batched TSAM Mask Generation}
\label{alg:mask-gen}
\begin{algorithmic}[1]
\REQUIRE Block tables $\mathbf{B}[B, L]$, metadata $\tau[N], t[N]$, query tenants $\mathbf{q}[B]$
\ENSURE Boolean mask $M[B, S_{\max}]$
\STATE \textbf{parallel for} $b \in [B]$, $\ell \in [L]$:
\STATE \quad $\mathrm{blk} \gets \mathbf{B}[b, \ell]$
\STATE \quad $M[b, \ell] \gets (\tau[\mathrm{blk}] = \textsc{Shared}) \;\lor\; (t[\mathrm{blk}] = \mathbf{q}[b])$
\end{algorithmic}
\end{algorithm}

\subsection{Flash Attention Integration}
\label{sec:design:flash}

\paragraph{Triton kernel.} Six additional lines in the inner loop of Flash Attention~\citep{dao2022flashattention, dao2023flashattention2, tillet2019triton}:
\begin{lstlisting}[language=Python, caption={TSAM in Triton Flash Attention.}]
mask_ptrs = Mask_ptr + bh * stride_mask_b + offs_n
tsam_mask = tl.load(mask_ptrs,
    mask=offs_n < seq_len_kv, other=False)
tsam_mask_2d = tsam_mask[None, :]
scores = tl.where(tsam_mask_2d, scores,
    float("-inf"))
\end{lstlisting}
Fully compatible with online softmax: $\exp(-\infty) = 0$ contributes nothing to running statistics.

\paragraph{vLLM CUDA kernel.} A single warp-uniform branch per block,
positioned after \texttt{qk}'s declaration so the variable is in scope:
\begin{lstlisting}[language=C++, caption={TSAM in vLLM CUDA kernel.}]
float qk = 0.0f;
if (page_type[phys_block] == PRIVATE
    && tenant_id[phys_block] != query_tenant) {
  qk = -INFINITY;  // IEEE 754: exp(-inf) = +0 exactly, matches reference
  continue;        // skip entire block
}
\end{lstlisting}
We use \texttt{-INFINITY} rather than \texttt{-FLT\_MAX} so that
$\exp(\mathrm{qk}) = +0$ \emph{exactly} per IEEE~754 \S6.1, matching the
Python reference (\texttt{torch.where(mask, scores, float("-inf"))}); with
\texttt{-FLT\_MAX} the resulting exponential is denormal but not exactly
zero, which would weaken the structural-zero argument.  The branch is
warp-uniform (all threads in a warp see the same page metadata) so there is
no divergence penalty.  Caller invariant: each query row must contain at
least one unmasked key, otherwise the softmax denominator is zero; the
calling Python wrapper enforces this via \texttt{assert\_any\_unmasked}
(see \texttt{tsam.integration.vllm\_patch}).  Reference-implementation
overhead: $\approx$3 instructions per block, amortized over
$16 \times 128$ multiply-accumulates ($\approx 0.015\%$ of dot-product
compute, before any kernel-fusion optimisation).

% =============================================================================
% Section 5: Formal Analysis
% =============================================================================

\section{Formal Analysis}
\label{sec:formal}

We now establish the theoretical foundations of TSAM, proving that typed
page metadata combined with attention masking provides \emph{exact
attention-channel isolation under correct page typing}: a structural
guarantee that the attention output at any tenant's positions is a
deterministic function of only the shared and own-private KV pairs.  We
state explicitly the assumptions under which this holds; these assumptions
are the deployment checklist that production operators must satisfy.

% -----------------------------------------------------------------------------
\subsection{Notation and Setup}
\label{sec:formal:notation}

Let $\mathcal{T} = \{t_1, t_2, \ldots, t_N\}$ denote the set of tenants concurrently
served within a single inference batch. Each \emph{typed page} is a tuple
$p = (\texttt{block\_id}, \tau, \rho, A)$, where:
\begin{itemize}[leftmargin=1.5em,itemsep=2pt]
    \item $\texttt{block\_id} \in \mathbb{N}$ is the unique physical page identifier in
          the KV cache;
    \item $\tau \in \{\textsc{System}, \textsc{Shared}, \textsc{Private}\}$ is the
          \emph{page type}, assigned at allocation time;
    \item $\rho \in \mathcal{T} \cup \{\bot\}$ is the \emph{owner field}, with
          $\rho = \bot$ for shared and system pages;
    \item $A \subseteq \mathcal{T}$ is the \emph{access set}, specifying which tenants
          may attend to this page.
\end{itemize}

\noindent
The following \textbf{type invariants} hold by construction:
\begin{align}
    \tau = \textsc{Shared}  &\implies \rho = \bot \;\wedge\; A = \mathcal{T}, \label{eq:inv-shared}\\
    \tau = \textsc{Private} &\implies \rho = t_i \;\wedge\; A = \{t_i\}. \label{eq:inv-private}
\end{align}

For a given tenant $t_i$, we define the \textbf{TSAM attention mask} $M_{t_i} \in \{0,1\}^L$
over the $L$ positions in the KV cache as:
\begin{equation}
    M_{t_i}[j] = \begin{cases}
        1 & \text{if } t_i \in A(p_j), \\
        0 & \text{otherwise},
    \end{cases}
    \label{eq:mask}
\end{equation}
where $p_j$ is the typed page containing position $j$, and $A(p_j)$ is its access set.

This induces a natural \textbf{three-way partition} of the KV positions for tenant $t_i$:
\begin{align}
    S      &= \{j : \tau(p_j) = \textsc{Shared}\},   && \text{(shared prefix)} \\
    O_i    &= \{j : \tau(p_j) = \textsc{Private} \wedge \rho(p_j) = t_i\}, && \text{(own private)} \\
    X_i    &= \{j : \tau(p_j) = \textsc{Private} \wedge \rho(p_j) \neq t_i\}. && \text{(cross-tenant private)}
\end{align}
By construction, $S \cup O_i \cup X_i = \{1, \ldots, L\}$ and the sets are pairwise disjoint.
From Equations~\eqref{eq:inv-shared}--\eqref{eq:mask}, we have $M_{t_i}[j] = 1$ for
$j \in S \cup O_i$ and $M_{t_i}[j] = 0$ for $j \in X_i$.

% -----------------------------------------------------------------------------
\subsection{Information-Theoretic Isolation Guarantee}
\label{sec:formal:isolation}

We establish perfect tenant isolation through a sequence of three results, each building
on the previous.

\begin{lemma}[Zero Attention Weight]
\label{lem:zero-weight}
For any tenant $t_i$ and any KV position $j \in X_i$, the attention weight
$w_{t_i}[j] = 0$ exactly.
\end{lemma}

\begin{proof}
The masked attention score at position $j$ is computed as:
\begin{equation}
    s_{t_i}[j] = \begin{cases}
        q_{t_i}^\top k_j / \sqrt{d_k} & \text{if } M_{t_i}[j] = 1, \\
        -\infty & \text{if } M_{t_i}[j] = 0.
    \end{cases}
\end{equation}
For $j \in X_i$, we have $M_{t_i}[j] = 0$, so $s_{t_i}[j] = -\infty$.
The softmax weight is:
\begin{equation}
    w_{t_i}[j] = \frac{\exp(s_{t_i}[j])}{Z} = \frac{\exp(-\infty)}{Z} = \frac{0}{Z} = 0,
\end{equation}
where $Z = \sum_{\ell=1}^{L} \exp(s_{t_i}[\ell]) > 0$ since at least one position
in $S \cup O_i$ has a finite score.

\medskip
\noindent\textbf{IEEE 754 exactness.}
Under IEEE 754 double-precision arithmetic, the mask value $-\infty$ is represented as
\texttt{-inf}, and $\exp(\texttt{-inf}) = +0$ \emph{exactly} (IEEE 754 \S6.1), not merely
a subnormal approximation. Division of exact zero by any finite positive $Z$ yields exact
zero. No rounding, underflow, or denormalization occurs---the zero is \emph{bitwise exact}.
In FlashAttention~\cite{dao2022flashattention}, the equivalent operation sets the
pre-softmax logit to \texttt{-inf} within each tile, producing the same exact zero
after the online softmax rescaling.
\end{proof}

\begin{lemma}[Functional Independence]
\label{lem:functional-independence}
Let $q$ be a query vector for tenant $t_i$, and let $K, V$ denote the full KV cache.
The attention output is a function only of the query and the visible KV pairs:
\begin{equation}
    \mathrm{Attn}_{t_i}(q, K, V) = f\!\big(q,\; \{(k_j, v_j)\}_{j \in S \cup O_i}\big).
\end{equation}
In particular, the output does not depend on any $(k_j, v_j)$ for $j \in X_i$.
\end{lemma}

\begin{proof}
The attention output for tenant $t_i$ decomposes as:
\begin{equation}
    \mathrm{Attn}_{t_i}(q, K, V) = \sum_{j=1}^{L} w_{t_i}[j] \cdot v_j
    = \underbrace{\sum_{j \in S \cup O_i} w_{t_i}[j] \cdot v_j}_{\text{visible contribution}}
    + \underbrace{\sum_{j \in X_i} w_{t_i}[j] \cdot v_j}_{\text{cross-tenant contribution}}.
\end{equation}
By Lemma~\ref{lem:zero-weight}, $w_{t_i}[j] = 0$ for all $j \in X_i$, so the
cross-tenant contribution vanishes identically:
\begin{equation}
    \sum_{j \in X_i} w_{t_i}[j] \cdot v_j = \sum_{j \in X_i} 0 \cdot v_j = \mathbf{0}.
\end{equation}
Furthermore, the normalization constant $Z$ also excludes cross-tenant positions:
\begin{equation}
    Z = \sum_{\ell=1}^{L} \exp(s_{t_i}[\ell])
      = \sum_{\ell \in S \cup O_i} \exp(s_{t_i}[\ell])
      + \sum_{\ell \in X_i} \underbrace{\exp(-\infty)}_{=\;0} = \sum_{\ell \in S \cup O_i} \exp(s_{t_i}[\ell]).
\end{equation}
Therefore, the weights $w_{t_i}[j]$ for $j \in S \cup O_i$ depend only on
$\{q, k_j : j \in S \cup O_i\}$, and the full output depends only on
$\{q, (k_j, v_j) : j \in S \cup O_i\}$.
\end{proof}

\begin{theorem}[Attention-Channel Noninterference under Correct Page Typing]
\label{thm:isolation}
Suppose the type invariants \eqref{eq:inv-shared}--\eqref{eq:inv-private}
hold and that the attention scores are computed exactly as in
Equation~\eqref{eq:mask}, with the mask applied before softmax using
$-\infty$ as the masking sentinel.
Then for any two distinct tenants $t_i, t_j \in \mathcal{T}$ with $i \neq j$,
the mutual information between tenant $t_i$'s attention output and tenant
$t_j$'s private KV data is exactly zero:
\begin{equation}
    I\!\big(\mathrm{Attn}_{t_i}(q, K, V);\; \mathrm{KV}_{\mathrm{private}}(t_j)\big) = 0.
\end{equation}
The result is a property of the \emph{attention channel} only: it does not
imply zero leakage through GPU memory side channels, model-mediated
semantic channels, allocation-pattern metadata, or any path other than
attention weights.  See Section~\ref{sec:limitations} for the full scope
boundary.
\end{theorem}

\begin{proof}
Define the following random variables:
\begin{align}
    A &\triangleq \mathrm{Attn}_{t_i}(q, K, V), &&\text{(tenant $t_i$'s attention output)} \\
    B &\triangleq \mathrm{KV}_{\mathrm{private}}(t_j) = \{(k_\ell, v_\ell)\}_{\ell \in O_j}, &&\text{(tenant $t_j$'s private KV pairs)} \\
    C &\triangleq \big(q,\; \{(k_\ell, v_\ell)\}_{\ell \in S \cup O_i}\big). &&\text{(tenant $t_i$'s visible context)}
\end{align}
Note that $O_j \subseteq X_i$ for $j \neq i$, so $B$ consists entirely of cross-tenant data.

\medskip
\noindent\textbf{Step 1: $A$ is a deterministic function of $C$.}
By Lemma~\ref{lem:functional-independence}, $A = f(C)$ for a deterministic function $f$.
This is a \emph{functional} relationship, not merely a statistical one.

\medskip
\noindent\textbf{Step 2: Conditional independence $A \perp\!\!\!\perp B \mid C$.}
Since $A = f(C)$ is completely determined by $C$, we have:
\begin{equation}
    P(A \mid C, B) = P(A \mid C) = \mathbf{1}_{A = f(C)},
\end{equation}
which is the definition of conditional independence $A \perp\!\!\!\perp B \mid C$.
The conditional mutual information satisfies:
\begin{equation}
    I(A; B \mid C) = H(A \mid C) - H(A \mid C, B) = 0 - 0 = 0,
    \label{eq:cond-mi-zero}
\end{equation}
since $A$ is deterministic given $C$.

\medskip
\noindent\textbf{Step 3: Statistical independence $C \perp\!\!\!\perp B$.}
In a multi-tenant serving scenario, each tenant's private input is generated independently:
tenant $t_j$'s private prompt (and hence KV pairs) is statistically independent of
tenant $t_i$'s private prompt and the shared prefix content, which together determine $C$.
Therefore:
\begin{equation}
    I(C; B) = 0.
    \label{eq:input-independence}
\end{equation}

\medskip
\noindent\textbf{Step 4: Contraction to zero.}
By the chain rule of mutual information~\cite{cover2006elements, shannon1948mathematical}:
\begin{equation}
    I(A; B) \leq I(A; B \mid C) + I(C; B).
\end{equation}
Substituting Equations~\eqref{eq:cond-mi-zero} and~\eqref{eq:input-independence}:
\begin{equation}
    I(A; B) \leq 0 + 0 = 0.
\end{equation}
Since mutual information is non-negative by definition~\cite{kraskov2004estimating}, we
conclude $I(A; B) = 0$, completing the proof.
\end{proof}

\begin{remark}[Comparison with Differential Privacy]
\label{rem:dp-comparison}
Theorem~\ref{thm:isolation} establishes \emph{exact} zero mutual information, which is
strictly stronger than $(\varepsilon, \delta)$-differential privacy for any
$\varepsilon > 0$, $\delta > 0$. Under differential privacy, an adversary retains
$\varepsilon$ nats of leakage per query, requiring careful privacy budget accounting
across interactions. TSAM's guarantee holds \emph{per query}, requires no privacy budget,
admits no degradation over time, and introduces no noise into the model output. This
corresponds to \emph{perfect secrecy} in the sense of Goguen and Meseguer's
non-interference~\cite{goguen1982security}: tenant $t_j$'s private inputs have
\emph{literally no effect} on tenant $t_i$'s outputs.
\end{remark}

\begin{remark}[Multi-Layer Extension]
\label{rem:multilayer}
For a transformer with $\mathcal{L}$ layers, the isolation guarantee extends by induction.
\textbf{Base case:} At layer $\ell = 1$, the KV cache entries are computed from embeddings
of each tenant's own tokens (private) or the shared prefix (shared); TSAM masking yields
$I(A^{(1)}_{t_i}; B) = 0$ by Theorem~\ref{thm:isolation}. \textbf{Inductive step:}
At layer $\ell + 1$, the input to attention is the output of layer $\ell$, which by the
inductive hypothesis depends only on $t_i$'s visible context. The typed page metadata is
preserved across layers (it is a property of the physical page, not the layer), so the
TSAM mask at layer $\ell + 1$ again zeros out all cross-tenant positions. Hence
$I(A^{(\ell+1)}_{t_i}; B) = 0$, and the result holds for all $\mathcal{L}$ layers.
\end{remark}

% -----------------------------------------------------------------------------
\subsection{Shared Prefix Safety}
\label{sec:formal:shared}

A natural concern is whether the shared prefix pages create an implicit information channel
between tenants. We show they do not.

\begin{theorem}[Shared Prefix Safety]
\label{thm:shared-safety}
Shared pages cannot serve as a channel for private data leakage between tenants. Formally,
for all tenants $t_i, t_j$ with $i \neq j$:
\begin{equation}
    I\!\big(\{(k_\ell, v_\ell)\}_{\ell \in S};\; \mathrm{KV}_{\mathrm{private}}(t_j)\big) = 0.
\end{equation}
\end{theorem}

\begin{proof}
The proof proceeds in three steps.

\medskip
\noindent\textbf{Step 1: Content invariance.}
Shared pages contain only the system prompt and shared prefix tokens, which are
\emph{identical} across all tenants and fixed before any private input is processed.
Let $\mathbf{x}_S$ denote the shared token sequence. The KV entries for shared pages are
$k_\ell = W_K \mathbf{x}_{S,\ell}$ and $v_\ell = W_V \mathbf{x}_{S,\ell}$, which are
deterministic functions of $\mathbf{x}_S$ and the (fixed) model weights. Since
$\mathbf{x}_S$ is independent of any tenant's private input, the shared KV content carries
no information about private data:
\begin{equation}
    I\!\big(\{(k_\ell, v_\ell)\}_{\ell \in S};\; \mathrm{KV}_{\mathrm{private}}(t_j)\big) = 0.
\end{equation}

\medskip
\noindent\textbf{Step 2: Uniform visibility.}
By invariant~\eqref{eq:inv-shared}, the access set for all shared pages is
$A = \mathcal{T}$. Every tenant observes \emph{the same} shared KV content with
\emph{the same} mask value ($M_{t_i}[j] = 1$ for all $i$ and all $j \in S$). There is no
tenant-specific variation in what is exposed, eliminating any covert signaling channel
through selective visibility.

\medskip
\noindent\textbf{Step 3: No write-back channel.}
The attention mechanism is a \emph{read-only} operation over the KV cache during the
decode phase: $\mathrm{Attn}(q, K, V) = \mathrm{softmax}(qK^\top/\sqrt{d_k})V$.
Tenant $t_i$'s queries do not modify the shared KV entries. While new tokens generated by
$t_i$ are appended to $O_i$ (private pages), they are never written to shared pages.
Therefore, tenant $t_i$ cannot encode information into the shared prefix that tenant $t_j$
could later read.
\end{proof}

% -----------------------------------------------------------------------------
\subsection{Scalability Analysis}
\label{sec:formal:scalability}

\begin{proposition}[TSAM Complexity]
\label{prop:complexity}
TSAM introduces the following overhead:
\begin{enumerate}[leftmargin=1.5em,itemsep=2pt]
    \item \textbf{Per-page metadata:} $O(1)$. Each typed page stores
          $(\tau, \rho) = (1~\text{byte} + 4~\text{bytes}) = 5$~bytes of metadata,
          independent of the number of tenants, sequence length, or model dimension.
    \item \textbf{Total metadata:} $O(N_{\mathrm{blocks}})$, where $N_{\mathrm{blocks}}$
          is the number of physical pages in the KV cache. Critically, this scales with
          \emph{memory capacity}, not with the number of tenants.
    \item \textbf{Mask generation:} $O(L)$ per query token, where $L$ is the total
          sequence length. Generating $M_{t_i}$ requires one table lookup per KV position
          to resolve the page type and compare the owner field.
    \item \textbf{Tenant scalability:} \emph{Unlimited}. The number of tenants $N$
          does not appear in any algebraic constraint on memory or compute. Adding a new
          tenant requires only allocating private pages and setting $\rho = t_{\mathrm{new}}$;
          no existing data structures are resized or reorganized.
\end{enumerate}
\end{proposition}

\begin{proof}
Claims (1)--(3) follow directly from the data structure definition in
\S\ref{sec:formal:notation}: the metadata is constant per page, total metadata is a linear
scan over pages, and mask generation requires one comparison per position. For claim (4),
observe that the mask condition $t_i \in A(p_j)$ reduces to checking $\rho(p_j) \in \{t_i, \bot\}$
under invariants~\eqref{eq:inv-shared}--\eqref{eq:inv-private}, which is $O(1)$ regardless
of $N$.
\end{proof}

% -----------------------------------------------------------------------------
\subsection{Comparison with Orthogonal Projection Methods}
\label{sec:formal:orthogonal}

Orthogonal projection methods~\cite{li2024orthogonal} achieve tenant isolation by mapping
each tenant's queries and keys into orthogonal subspaces. We show this approach suffers
from fundamental limitations that TSAM avoids.

\begin{definition}[Orthogonal Projection Isolation]
Each tenant $t_i$ is assigned a projection matrix $P_i = U_i U_i^\top \in \mathbb{R}^{d_{\mathrm{model}} \times d_{\mathrm{model}}}$,
where $U_i \in \mathbb{R}^{d_{\mathrm{model}} \times d_{\mathrm{sub}}}$ has orthonormal columns
spanning a $d_{\mathrm{sub}}$-dimensional subspace. Isolation requires mutual orthogonality:
$P_i P_j = 0$ for all $i \neq j$.
\end{definition}

\begin{proposition}[Tenant Capacity Bound]
\label{prop:ortho-capacity}
Orthogonal projection isolation admits at most
$N_{\max} = \lfloor d_{\mathrm{model}} / d_{\mathrm{sub}} \rfloor$
tenants. For typical configurations ($d_{\mathrm{model}} = 4096$, $d_{\mathrm{sub}} = 256$),
this yields $N_{\max} = 16$.
\end{proposition}

\begin{proof}
The orthogonality constraint $P_i P_j = 0$ requires $\mathrm{range}(U_i) \perp \mathrm{range}(U_j)$
for all $i \neq j$. Since $\mathbb{R}^{d_{\mathrm{model}}}$ can be decomposed into at most
$\lfloor d_{\mathrm{model}} / d_{\mathrm{sub}} \rfloor$ mutually orthogonal subspaces
of dimension $d_{\mathrm{sub}}$, the bound follows. With $d_{\mathrm{model}} = 4096$ and
$d_{\mathrm{sub}} = 256$: $N_{\max} = \lfloor 4096 / 256 \rfloor = 16$.
\end{proof}

\begin{proposition}[Prefix Sharing Incompatibility]
\label{prop:ortho-prefix}
Orthogonal projection fundamentally prevents shared prefix optimization. For a shared
key vector $k \in \mathbb{R}^{d_{\mathrm{model}}}$, the projected keys seen by different
tenants differ:
\begin{equation}
    P_i k \neq P_j k \quad \text{for generic } k \text{ and } i \neq j,
\end{equation}
since $P_i$ and $P_j$ project onto orthogonal subspaces. Consequently, each tenant requires
a separate copy of the shared prefix in its projected subspace, eliminating the
$O(N)$-to-$O(1)$ memory savings that prefix sharing provides.
\end{proposition}

\begin{proof}
Let $k = k_\parallel^{(i)} + k_\perp^{(i)}$ be the decomposition with respect to
$\mathrm{range}(U_i)$. Then $P_i k = k_\parallel^{(i)}$ and $P_j k = k_\parallel^{(j)}$.
Since $\mathrm{range}(U_i) \perp \mathrm{range}(U_j)$, we have
$\langle P_i k, P_j k \rangle = 0$, so $P_i k = P_j k$ only if both projections are zero.
For a generic $k$ (i.e., $k$ not in the null space of both $P_i$ and $P_j$), the projected
keys are distinct nonzero vectors in orthogonal subspaces. Therefore, the physical KV cache
entries cannot be shared---each tenant requires its own projected copy.
\end{proof}

\begin{table}[t]
\centering
\caption{Comparison of TSAM and orthogonal projection isolation. TSAM achieves
strictly superior guarantees across all axes.}
\label{tab:comparison}
\small
\begin{tabular}{@{}lcc@{}}
\toprule
\textbf{Property} & \textbf{Orthogonal Projection} & \textbf{TSAM} \\
\midrule
Isolation guarantee & $I(A;B) = 0$ & $I(A;B) = 0$ \\
Max tenants & $\lfloor d_{\mathrm{model}}/d_{\mathrm{sub}} \rfloor$ (e.g., 16) & \emph{Unlimited} \\
Prefix sharing & \ding{55}~Incompatible & \ding{51}~Native support \\
Model modification & Required (projection layers) & None (mask only) \\
Representation capacity & Reduced ($d_{\mathrm{sub}} < d_{\mathrm{model}}$) & Full $d_{\mathrm{model}}$ \\
Per-tenant memory overhead & $O(d_{\mathrm{model}} \cdot d_{\mathrm{sub}})$ & $O(1)$ (5 bytes/page) \\
\bottomrule
\end{tabular}
\end{table}

% =============================================================================
% Section 6: End-to-End Noninterference Across the Full Transformer Stack
% =============================================================================

\section{End-to-End Noninterference Across the Full Transformer Stack}
\label{sec:e2e}

The Theorem~\ref{thm:isolation} establishes that a \emph{single} attention
layer with TSAM masking satisfies $I = 0$ between the attacker's attention
output and the victim's private KV entries.  A natural and critical
question follows:

\begin{center}
\emph{Does per-layer TSAM masking guarantee end-to-end noninterference\\
through an $L$-layer transformer with residual connections, LayerNorm, and MLPs?}
\end{center}

We answer affirmatively, under precisely characterized conditions.

% -----------------------------------------------------------------------------
\subsection{Sufficient Conditions for End-to-End Isolation}
\label{sec:e2e:conditions}

\begin{definition}[TSAM-Compliant Transformer]
\label{def:compliant}
An $L$-layer decoder-only transformer is \textbf{TSAM-compliant} if:
\begin{enumerate}[label=\textbf{C\arabic*}.,leftmargin=2em]
    \item \textbf{Per-layer masking.}  Every attention sublayer applies a TSAM mask
          $M_{t_i}$ that blocks all cross-tenant private positions ($j \in X_i$).
    \item \textbf{Clean shared prefix.}  All SHARED KV-cache entries are computed
          from public inputs only (system prompt $+$ frozen model weights), with
          no private tenant data in the computation batch.
    \item \textbf{Position-local non-attention sublayers.}  LayerNorm, MLP, and activation
          functions operate independently on each sequence position (no cross-position
          mixing outside of attention).
    \item \textbf{Frozen shared weights.}  Model parameters $\theta$ are identical
          for all tenants (no per-tenant adapters within the masked computation).
\end{enumerate}
\end{definition}

\noindent
Condition \textbf{C3} is satisfied by construction in standard transformer
architectures (GPT, LLaMA, Qwen, Gemma): LayerNorm normalises along the
feature dimension per position, and MLPs apply the same pointwise transform
at each position independently.  Condition~\textbf{C2}---a clean shared
prefix---is the most easily violated in production: vLLM's prefix caching
computes shared KV entries during batches that may include private tokens,
so deployments must enforce a clean-compute provenance check before
admitting a page as SHARED.  Conditions~\textbf{C1} and~\textbf{C4} are
operator-controlled deployment choices; \S\ref{sec:limitations} discusses
what happens when each is violated.

% -----------------------------------------------------------------------------
\subsection{Formal Guarantee}
\label{sec:e2e:theorem}

\begin{theorem}[End-to-End Noninterference under C1--C4]
\label{thm:e2e}
Let $\mathcal{M}$ be a TSAM-compliant $L$-layer transformer
(Definition~\ref{def:compliant}).  For any two distinct tenants
$t_i, t_j$ with $i \neq j$, the mutual information between tenant
$t_i$'s output at \emph{any} layer $\ell \in \{1, \ldots, L\}$
and tenant $t_j$'s private input is exactly zero:
\begin{equation}
    I\!\big(h^{(\ell)}_{t_i};\; \mathrm{tokens}_{\mathrm{private}}(t_j)\big) = 0
    \quad \forall\, \ell \in \{1, \ldots, L\}.
\end{equation}
In particular, the final output logits satisfy
$I\!\big(\mathrm{logits}_{t_i};\; \mathrm{tokens}_{\mathrm{private}}(t_j)\big) = 0$.
\end{theorem}

\begin{proof}
By structural induction on the layer index $\ell$.

\medskip\noindent
\textbf{Base case ($\ell = 0$, embedding layer).}
The input to the first transformer layer is
$h^{(0)}[p] = \mathrm{Embed}(\mathrm{token}[p]) + \mathrm{PosEmbed}(p)$.
At position $p$ belonging to tenant $t_i$, the embedding depends only on
the token at position $p$ and the position index.  Since tenant $t_j$'s
private tokens occupy different positions, $h^{(0)}$ at $t_i$'s positions
is structurally independent of $t_j$'s tokens.

\medskip\noindent
\textbf{Inductive step ($\ell - 1 \to \ell$).}
Assume $h^{(\ell-1)}$ at all of $t_i$'s positions is independent of
$\mathrm{tokens}_{\mathrm{private}}(t_j)$.  Layer $\ell$ computes:
\begin{align}
    a^{(\ell)}[p] &= h^{(\ell-1)}[p] + \mathrm{Attn}^{(\ell)}\!\big(\mathrm{LN}(h^{(\ell-1)}),\; M_{t_i}\big)[p], \label{eq:attn-sublayer}\\
    h^{(\ell)}[p] &= a^{(\ell)}[p] + \mathrm{MLP}^{(\ell)}\!\big(\mathrm{LN}(a^{(\ell)}[p])\big). \label{eq:mlp-sublayer}
\end{align}

\noindent\emph{Attention sublayer} (Eq.~\ref{eq:attn-sublayer}, by C1):
By Theorem~\ref{thm:isolation}, $\mathrm{Attn}^{(\ell)}$ with mask $M_{t_i}$
produces output that depends only on positions in $S \cup O_i$.  The keys
and values at these positions are computed from $h^{(\ell-1)}$ at
$S \cup O_i$, which by the inductive hypothesis (and by C2 for $S$) are
independent of $t_j$'s private data.  Therefore
$\mathrm{Attn}^{(\ell)}[p] \perp\!\!\!\perp \mathrm{tokens}_{\mathrm{private}}(t_j)$.

\noindent\emph{Residual connection}: Adding $h^{(\ell-1)}[p]$ (independent
by IH) to $\mathrm{Attn}^{(\ell)}[p]$ (independent by above) preserves
independence.

\noindent\emph{LayerNorm and MLP} (by C3): Operate per-position along the
feature dimension; no cross-position mixing therefore no cross-tenant
information flow.  The MLP's pointwise transformation, applied to an
independent input, produces an independent output.

\noindent\emph{Second residual}: Same argument.  Therefore $h^{(\ell)}[p]$
is independent of $\mathrm{tokens}_{\mathrm{private}}(t_j)$.

\medskip\noindent
\textbf{Final output (by C4).}
The output projection
$\mathrm{logits} = W_{\mathrm{out}} \cdot \mathrm{LN}(h^{(L)})$
is a deterministic per-position linear map applied to $h^{(L)}$, with
$W_{\mathrm{out}}$ identical across tenants by C4.  By the induction
result at $\ell = L$, the logits at $t_i$'s positions are independent of
$t_j$'s private data, completing the proof.
\end{proof}

\begin{remark}[Tightness of Conditions C1--C4]
\label{rem:tightness}
Each condition is necessary; removing any one breaks the chain.
\textbf{C1}: a single unmasked layer admits cross-tenant attention.
\textbf{C2}: contaminated shared KV entries carry information from one
tenant's private data to all others.
\textbf{C3}: a cross-position operation (e.g., sequence-level
normalization, 1D convolution over positions) would mix tenant signals.
\textbf{C4}: per-tenant LoRA adapters in the masked attention path could
encode tenant-specific information into the shared computation.  These
four conditions are the deployment checklist for any production system
that wishes to invoke Theorem~\ref{thm:e2e}.
\end{remark}

% -----------------------------------------------------------------------------
\subsection{Empirical Verification}
\label{sec:e2e:empirical}

We verify Theorem~\ref{thm:e2e} empirically on a TSAM-compliant transformer
implementation up to 24 layers.  All experiments are reproducible from the
\texttt{tsam.evaluation.e2e\_noninterference} module in the open-source
release.

\paragraph{Experimental setup.}
For each depth $L \in \{1, 2, 4, 8, 12, 24\}$, we instantiate a transformer
($d_{\mathrm{model}} = 256$, $H = 4$ heads, $d_{\mathrm{ff}} = 1024$,
vocabulary size $1000$) with random weights.  We construct worst-case
batches:
$[\text{shared}_{16}\,|\,\text{victim}_{16}\,|\,\text{attacker}_{16}]$
where the victim's positions precede the attacker's, maximising the causal
attack surface.  For each trial, we run the model twice with different
victim tokens and measure the L2 norm of the difference in the attacker's
output logits.  Under TSAM, the victim positions are masked; under vanilla,
full causal attention is used.

\paragraph{Functional independence.}

\begin{table}[h]
\centering
\caption{Functional independence (FI): $\|f(x_A) - f(x_B)\|_2$ at attacker
positions, where $x_A, x_B$ differ only in the victim's private tokens.
50 trials per cell.  Measurements from
\texttt{tsam.evaluation.e2e\_noninterference}.}
\label{tab:fi}
\begin{tabular}{rcc}
\toprule
Depth $L$ & TSAM (max over 50 trials) & Vanilla (mean) \\
\midrule
1  & \textbf{0.0 (exact)} & 1.88 \\
2  & \textbf{0.0 (exact)} & 2.70 \\
4  & \textbf{0.0 (exact)} & 3.69 \\
8  & \textbf{0.0 (exact)} & 5.08 \\
12 & \textbf{0.0 (exact)} & 6.28 \\
24 & \textbf{0.0 (exact)} & 8.07 \\
\bottomrule
\end{tabular}
\end{table}

TSAM achieves bitwise-exact zero at every depth from 1 to 24 layers.
Vanilla leakage grows monotonically with depth as cross-tenant signal
accumulates through residual connections, consistent with prior empirical
observations on attention-driven information flow~\citep{carlini2021extracting,
nasr2023scalable}.

\paragraph{Linear probing attack.}
We train a linear classifier on the attacker's output logits to predict
which of $K = 8$ victim private prompts was used.  Random baseline
$1/8 = 12.5\%$.

\begin{table}[h]
\centering
\caption{Linear probing accuracy (\%) predicting the victim's private prompt
class from the attacker's output logits.  200 samples, 70/30 train/test
split, ridge regression.}
\label{tab:probe}
\begin{tabular}{rcc}
\toprule
Depth $L$ & TSAM & Vanilla \\
\midrule
1  & 11.7\% ($\approx$ random) & \textbf{100\%} \\
4  & 11.7\% ($\approx$ random) & \textbf{100\%} \\
12 & 11.7\% ($\approx$ random) & \textbf{100\%} \\
24 & 11.7\% ($\approx$ random) & \textbf{100\%} \\
\bottomrule
\end{tabular}
\end{table}

Under TSAM, the probe achieves random chance at every depth.  Under vanilla
attention, a simple linear classifier achieves perfect extraction at every
depth---demonstrating that, without isolation, cross-tenant information
flows freely through unmasked transformers and is trivially recoverable.

\paragraph{Key finding.}
Per-layer TSAM attention masking is sufficient for end-to-end noninterference
through the full transformer stack, provided conditions C1--C4 hold.  The
isolation is not approximate, not statistical, and does not degrade with
depth---it is a structural invariant preserved by induction through all
$L$ layers.

% =============================================================================
% Section 7: DAI: Timing Side-Channel Defense
% =============================================================================

\section{DAI: Timing Side-Channel Defense}
\label{sec:dai}

Theorem~\ref{thm:isolation} guarantees that no private data leaks through the attention
\emph{data channel}. However, the shared KV cache introduces a potential \emph{metadata
channel}: cache hit/miss timing may reveal whether a page is resident, leaking occupancy
patterns~\cite{osvik2006cache, yarom2014flush, percival2005cache}. We introduce the
\textbf{Deterministic Anomaly Interceptor (DAI)}, a lightweight per-tenant state machine
that detects and mitigates timing side-channel attacks in real time.

% -----------------------------------------------------------------------------
\subsection{Collision Attack Model}
\label{sec:dai:attack}

We adopt the standard cache timing oracle model~\cite{osvik2006cache, percival2005cache}.
An adversary controlling tenant $t_a$ issues carefully crafted queries to probe whether
a target tenant $t_v$'s pages are resident in the KV cache. The observed latency for a
single probe is:
\begin{equation}
    T_{\mathrm{obs}} = \begin{cases}
        T_{\mathrm{hit}} + \eta & \text{if the target page is cached}, \\
        T_{\mathrm{miss}} + \eta & \text{if the target page is evicted},
    \end{cases}
    \label{eq:timing-oracle}
\end{equation}
where $T_{\mathrm{hit}} < T_{\mathrm{miss}}$ are the deterministic cache hit and miss
latencies, $\eta \sim \mathcal{N}(0, \sigma_{\mathrm{sys}}^2)$ is system noise, and
$\Delta \triangleq T_{\mathrm{miss}} - T_{\mathrm{hit}}$ is the timing gap.

Let $S_{\mathrm{cache}} \in \{0, 1\}$ denote the binary cache state (0 = miss, 1 = hit).
The adversary's observation forms a \emph{binary channel}~\cite{shannon1948mathematical}
with capacity:
\begin{equation}
    C = 1 - H_b(p_e),
    \label{eq:channel-capacity}
\end{equation}
where $H_b(p_e) = -p_e \log_2 p_e - (1 - p_e)\log_2(1 - p_e)$ is the binary entropy
function and $p_e = \Phi(-\Delta / (2\sigma_{\mathrm{sys}}))$ is the crossover probability
under optimal thresholding ($\Phi$ is the standard normal CDF).

\textbf{Without mitigation}, the adversary can average over $n$ independent probes,
reducing the effective noise to $\sigma_{\mathrm{eff}} = \sigma_{\mathrm{sys}} / \sqrt{n}$.
As $n \to \infty$, $p_e \to 0$, yielding $C \to 1$ bit per probe---a perfect timing oracle.

% -----------------------------------------------------------------------------
\subsection{DFA-Based Detection}
\label{sec:dai:dfa}

DAI models each tenant's behavior as a \textbf{deterministic finite automaton (DFA)} with
four states, enabling stateful detection that distinguishes transient anomalies from
sustained attack patterns.

\begin{definition}[DAI State Machine]
The DFA is the tuple $(\mathcal{S}, \Sigma, \delta, s_0, \mathcal{F})$ where:
\begin{itemize}[leftmargin=1.5em,itemsep=2pt]
    \item $\mathcal{S} = \{S_0, S_1, S_2, S_3\}$: states \textsc{Normal},
          \textsc{Probe\_1}, \textsc{Probe\_2}, \textsc{Attack\_Confirmed};
    \item $\Sigma$: per-query statistics (rolling variance $\hat{\sigma}$, query count,
          distributional divergence);
    \item $\delta: \mathcal{S} \times \Sigma \to \mathcal{S}$: transition function
          (Table~\ref{tab:dfa-transitions});
    \item $s_0 = S_0$: initial state;
    \item $\mathcal{F} = \{S_3\}$: the accepting (defense-active) state.
\end{itemize}
\end{definition}

\noindent The transition guards are parameterized by three thresholds:

\begin{enumerate}[leftmargin=1.5em,itemsep=2pt]
    \item \textbf{$S_0 \to S_1$ (anomaly detection):} $\hat{\sigma} > \tau_\sigma$, where
          $\hat{\sigma}$ is the rolling standard deviation of inter-query arrival times
          over a sliding window of $W = 50$ queries, and $\tau_\sigma = 15$~ms.
    \item \textbf{$S_1 \to S_2$ (sustained anomaly):} the condition $\hat{\sigma} > \tau_\sigma$
          persists for $K = 20$ consecutive queries, ruling out transient load spikes.
    \item \textbf{$S_2 \to S_3$ (distributional confirmation):} the KL divergence between
          the observed timing distribution $\hat{P}$ and the benign baseline $P_{\mathrm{benign}}$
          exceeds the threshold:
          $D_{\mathrm{KL}}(\hat{P} \| P_{\mathrm{benign}}) > \kappa$, with $\kappa = 2.0$~nats.
    \item \textbf{$S_3 \to S_0$ (cooldown):} the tenant must satisfy \emph{all three}
          conditions simultaneously for at least $C_{\mathrm{cool}} = 100$ consecutive
          queries: $\hat{\sigma} \leq \tau_\sigma$, $D_{\mathrm{KL}}(\hat{P} \| P_{\mathrm{benign}}) \leq \kappa$,
          and the query rate returns to within $2\times$ of the tenant's historical baseline.
\end{enumerate}

\begin{table}[t]
\centering
\caption{DAI state transition table. Each row specifies the current state, guard condition,
next state, and triggered action.}
\label{tab:dfa-transitions}
\small
\begin{tabular}{@{}ccll@{}}
\toprule
\textbf{From} & \textbf{To} & \textbf{Guard} & \textbf{Action} \\
\midrule
$S_0$ & $S_1$ & $\hat{\sigma} > \tau_\sigma$ & Begin monitoring window \\
$S_1$ & $S_0$ & $\hat{\sigma} \leq \tau_\sigma$ before $K$ queries & Reset counter \\
$S_1$ & $S_2$ & $\hat{\sigma} > \tau_\sigma$ sustained for $K = 20$ queries & Compute $D_{\mathrm{KL}}$ \\
$S_2$ & $S_0$ & $D_{\mathrm{KL}}(\hat{P} \| P_{\mathrm{benign}}) \leq \kappa$ & False alarm; reset \\
$S_2$ & $S_3$ & $D_{\mathrm{KL}}(\hat{P} \| P_{\mathrm{benign}}) > \kappa = 2.0$ nats & Activate countermeasures \\
$S_3$ & $S_3$ & Any guard violated during cooldown & Maintain defense \\
$S_3$ & $S_0$ & $C_{\mathrm{cool}} = 100$ clean queries & Deactivate; resume normal \\
\bottomrule
\end{tabular}
\end{table}

\begin{algorithm}[t]
\caption{DAI Per-Tenant State Machine}
\label{alg:dai}
\begin{algorithmic}[1]
\Require Tenant state $s \in \{S_0, S_1, S_2, S_3\}$, timing buffer $\mathcal{B}$, counter $c$
\Ensure Updated state $s'$, response latency modifier $\delta_t$
\State $\mathcal{B}.\text{append}(T_{\mathrm{current}})$ \Comment{Record query timestamp}
\State $\hat{\sigma} \gets \text{RollingStd}(\mathcal{B}, W{=}50)$
\If{$s = S_0$}
    \If{$\hat{\sigma} > \tau_\sigma$}
        $s' \gets S_1$; $c \gets 1$
    \Else~$s' \gets S_0$
    \EndIf
\ElsIf{$s = S_1$}
    \If{$\hat{\sigma} \leq \tau_\sigma$}
        $s' \gets S_0$; $c \gets 0$ \Comment{Transient; reset}
    \ElsIf{$c \geq K$}
        \State $\hat{P} \gets \text{HistogramEstimate}(\mathcal{B})$
        \If{$D_{\mathrm{KL}}(\hat{P} \| P_{\mathrm{benign}}) > \kappa$}
            $s' \gets S_3$ \Comment{Skip $S_2$: confirmed}
        \Else~$s' \gets S_0$; $c \gets 0$ \Comment{Benign anomaly}
        \EndIf
    \Else~$c \gets c + 1$; $s' \gets S_1$
    \EndIf
\ElsIf{$s = S_3$}
    \State $\delta_t \gets |\mathcal{N}(0, \sigma_{\mathrm{noise}}^2)|$ \Comment{Inject noise}
    \State Apply rate limiting: delay by factor $\lambda = 2.0$
    \If{$\hat{\sigma} \leq \tau_\sigma ~\wedge~ D_{\mathrm{KL}} \leq \kappa ~\wedge~ c_{\mathrm{cool}} \geq C_{\mathrm{cool}}$}
        $s' \gets S_0$ \Comment{Cooldown complete}
    \Else~$s' \gets S_3$; update $c_{\mathrm{cool}}$
    \EndIf
\EndIf
\State \Return $s', \delta_t$
\end{algorithmic}
\end{algorithm}

% -----------------------------------------------------------------------------
\subsection{Countermeasures}
\label{sec:dai:countermeasures}

Upon entering state $S_3$, DAI activates two complementary countermeasures that jointly
destroy the timing channel.

\paragraph{Noise injection.}
Each response to the flagged tenant is delayed by an additive noise term
$\delta_t = |\mathcal{N}(0, \sigma_{\mathrm{noise}}^2)|$ with
$\sigma_{\mathrm{noise}} = 5$~ms (half-normal distribution, ensuring non-negative delay).
This noise is \emph{independent} of the cache state $S_{\mathrm{cache}}$ and is injected
at the serving layer, after attention computation.

\paragraph{Rate limiting.}
The flagged tenant's query throughput is throttled by a factor of $\lambda = 2.0$,
doubling the minimum inter-query interval. This limits the number of samples available
for averaging attacks.

\paragraph{SNR analysis.}
The pre-defense signal-to-noise ratio for a single probe is:
\begin{equation}
    \mathrm{SNR}_{\mathrm{pre}} = \frac{\Delta^2}{\sigma_{\mathrm{sys}}^2}.
\end{equation}
After noise injection, the effective noise variance becomes
$\sigma_{\mathrm{eff}}^2 = \sigma_{\mathrm{sys}}^2 + \sigma_{\mathrm{noise}}^2$, yielding:
\begin{equation}
    \mathrm{SNR}_{\mathrm{post}} = \frac{\Delta^2}{\sigma_{\mathrm{sys}}^2 + \sigma_{\mathrm{noise}}^2}
    \ll \mathrm{SNR}_{\mathrm{pre}}.
    \label{eq:snr-post}
\end{equation}
For typical values ($\Delta = 2$~ms, $\sigma_{\mathrm{sys}} = 1$~ms,
$\sigma_{\mathrm{noise}} = 5$~ms), we have
$\mathrm{SNR}_{\mathrm{post}} = 4 / 26 \approx 0.15$, a $17\times$ reduction.

\paragraph{Effective channel capacity.}
Combining noise injection with rate limiting, the effective channel capacity becomes:
\begin{equation}
    C_{\mathrm{eff}} = \frac{1}{\lambda}\big(1 - H_b(p_e')\big),
    \label{eq:effective-capacity}
\end{equation}
where $p_e' = \Phi\!\big(-\Delta / (2\sqrt{\sigma_{\mathrm{sys}}^2 + \sigma_{\mathrm{noise}}^2})\big)$.
Substituting:
\begin{equation}
    p_e' = \Phi\!\bigg(\frac{-2}{2\sqrt{26}}\bigg) = \Phi(-0.196) \approx 0.422.
\end{equation}
The binary entropy evaluates to $H_b(0.422) \approx 0.982$ bits, giving:
\begin{equation}
    C_{\mathrm{eff}} = \frac{1}{2}(1 - 0.982) = 0.009~\text{bits/probe}.
\end{equation}
This represents a $>99\%$ reduction in channel capacity compared to the undefended
$C \to 1$ bit/probe. Furthermore, rate limiting halves the probe rate, so the
\emph{throughput} of leaked information drops by an additional factor of $2$, yielding
an effective information rate of $\sim$0.0045 bits per unit time---practically zero.

% -----------------------------------------------------------------------------
\subsection{Timing Information Bound}
\label{sec:dai:bound}

\begin{theorem}[Timing Information Bound]
\label{thm:timing-bound}
Under DAI, the mutual information between the cache state and observed timing is bounded:
\begin{equation}
    I(S_{\mathrm{cache}};\; T_{\mathrm{obs}}) \leq I_{\mathrm{natural}} + \varepsilon_{\mathrm{FNR}},
\end{equation}
where $I_{\mathrm{natural}}$ is the mutual information under benign operation
(due to inherent system timing variation) and $\varepsilon_{\mathrm{FNR}}$ is a residual
term bounded by the false negative rate of the DFA detector.
\end{theorem}

\begin{proof}
We analyze three temporal phases of any attack.

\medskip
\noindent\textbf{Phase 1: Pre-detection (states $S_0 \to S_1 \to S_2$).}
Before DAI transitions to $S_3$, the adversary operates with the undefended channel.
The maximum number of queries in this phase is bounded by the detection latency:
$n_{\mathrm{pre}} \leq K + W = 20 + 50 = 70$ queries (the monitoring window $W$ plus
the sustained anomaly threshold $K$). During this phase, the adversary accumulates at most:
\begin{equation}
    I_{\mathrm{pre}} \leq n_{\mathrm{pre}} \cdot C \leq 70 \cdot 1 = 70~\text{bits}
\end{equation}
of timing information. This is a \emph{finite, bounded} quantity, independent of the
total attack duration.

\medskip
\noindent\textbf{Phase 2: Active defense (state $S_3$).}
Once countermeasures activate, the observed timing becomes:
\begin{equation}
    T_{\mathrm{obs}} = T_{\mathrm{true}} + \delta_t, \quad \delta_t \sim |\mathcal{N}(0, \sigma_{\mathrm{noise}}^2)|,
    \quad \delta_t \perp\!\!\!\perp S_{\mathrm{cache}}.
\end{equation}
By the \textbf{data processing inequality}~\cite{cover2006elements}:
\begin{equation}
    I(S_{\mathrm{cache}}; T_{\mathrm{obs}}) \leq I(S_{\mathrm{cache}}; T_{\mathrm{true}}).
\end{equation}
But $T_{\mathrm{obs}}$ is a noisy function of $T_{\mathrm{true}}$, and the additive
independent noise strictly reduces mutual information. Specifically, by
Equation~\eqref{eq:effective-capacity}, the per-probe capacity drops to
$C_{\mathrm{eff}} \approx 0.009$ bits. Combined with rate limiting ($\lambda = 2.0$),
the information rate is:
\begin{equation}
    R_{\mathrm{attack}} = \frac{C_{\mathrm{eff}}}{\lambda} \approx 0.0045~\text{bits/query} \to 0
    \quad \text{as } \sigma_{\mathrm{noise}} / \Delta \to \infty.
\end{equation}

\medskip
\noindent\textbf{Phase 3: Steady state.}
Under a persistent attack ($\hat{\sigma}$ remains elevated), the DFA remains in $S_3$
indefinitely. The cooldown guard ($C_{\mathrm{cool}} = 100$ consecutive clean queries
with $\hat{\sigma} \leq \tau_\sigma$ \emph{and} $D_{\mathrm{KL}} \leq \kappa$) prevents
premature relaxation. If the attacker ceases probing to satisfy the cooldown condition,
they are by definition no longer attacking, and the channel is idle.

If the attacker resumes after cooldown, the DFA re-enters $S_1$ within one anomalous
query and reaches $S_3$ within $K = 20$ queries, bounding the new pre-detection leakage
to another $\leq 70$ bits---a pattern that yields at most $70$ bits per attack epoch,
with epoch length growing due to the mandatory $100$-query cooldown.

\medskip
\noindent\textbf{Combining the phases.}
The total steady-state leakage rate is:
\begin{equation}
    I(S_{\mathrm{cache}}; T_{\mathrm{obs}}) \leq I_{\mathrm{natural}} + \varepsilon_{\mathrm{FNR}},
\end{equation}
where $I_{\mathrm{natural}}$ captures the unavoidable timing variation of benign cache
behavior (present in \emph{any} caching system) and
$\varepsilon_{\mathrm{FNR}} = P(\text{FN}) \cdot C$ accounts for the residual leakage
during any undetected attack queries. In our empirical evaluation (\S\ref{sec:eval}),
the DFA achieves a false negative rate below $10^{-3}$, yielding
$\varepsilon_{\mathrm{FNR}} < 10^{-3}$ bits/query.
\end{proof}

\begin{remark}[Defense in Depth]
\label{rem:defense-in-depth}
TSAM and DAI provide complementary, layered protection~\cite{denning1976lattice, goguen1982security}:
\begin{enumerate}[leftmargin=1.5em,itemsep=2pt]
    \item \textbf{Data channel} (TSAM): Theorem~\ref{thm:isolation} guarantees
          $I(\mathrm{Attn}_{t_i}; \mathrm{KV}_{\mathrm{private}}(t_j)) = 0$ exactly---the
          attention mechanism reveals \emph{zero bits} of private data.
    \item \textbf{Metadata channel} (DAI): Theorem~\ref{thm:timing-bound} bounds the
          timing side-channel to $I_{\mathrm{natural}} + \varepsilon_{\mathrm{FNR}}
          \approx 0$---the cache timing reveals \emph{negligible bits} of occupancy
          metadata.
\end{enumerate}
Together, these two mechanisms close both the primary data exfiltration vector and the
secondary timing side-channel, achieving comprehensive tenant isolation without modifying
model weights, degrading output quality, or imposing significant computational overhead.
\end{remark}


% ==============================================================================
% SECTION 7 — IMPLEMENTATION
% ==============================================================================
\section{Implementation}
\label{sec:implementation}

\subsection{Core Type System}

TSAM enforces invariants through Python's type system. All metadata structures are frozen dataclasses with \texttt{\_\_post\_init\_\_} validation:

\begin{lstlisting}[language=Python, caption={Core type definitions.}]
class PageType(IntEnum):
    PRIVATE = 0
    SHARED  = 1

@dataclass(frozen=True)
class TypedPage:
    block_id: int
    tenant_id: int
    page_type: PageType
    access_list: Optional[FrozenSet[int]]

    def __post_init__(self):
        if self.page_type == PageType.PRIVATE:
            assert self.tenant_id != SHARED_TENANT_ID
            assert self.access_list == frozenset({self.tenant_id})
        else:
            assert self.tenant_id == SHARED_TENANT_ID
            assert self.access_list is None
\end{lstlisting}

The GPU-resident metadata tensors (\texttt{uint8} for page type, \texttt{int32} for tenant ID) avoid host--device transfers on the critical path. For 65,536 blocks, total metadata is 320\,KB---negligible.

\subsection{Triton Flash Attention Kernel}

The modified Flash Attention kernel~\citep{dao2022flashattention, tillet2019triton} uses tile sizes $\textsc{Block\_M} = 64$, $\textsc{Block\_N} = 64$ with online softmax. The TSAM mask is loaded per-tile and broadcast to $(64, 64)$; blocked positions receive $-\infty$ before the softmax accumulation. The kernel is fully compatible with GQA~\citep{ainslie2023gqa} and MQA~\citep{shazeer2019fast}: the mask operates on the KV-head level and broadcasts to query heads.

\subsection{vLLM Integration}

Three integration points:
\begin{enumerate}[leftmargin=*,itemsep=2pt]
    \item \textbf{TSAMBlockTableExtension}: Wraps vLLM's block manager with type metadata tensors. Provides \texttt{mark\_private}/\texttt{mark\_shared} and \texttt{generate\_mask} methods.
    \item \textbf{TSAMAttentionWrapper}: Intercepts the attention \texttt{forward()} call, constructs the TSAM mask, and combines it with existing attention masks via element-wise AND.
    \item \textbf{Block manager patch}: Monkey-patches \texttt{allocate()} to propagate tenant metadata during block allocation.
\end{enumerate}

The CUDA kernel patch adds a single warp-uniform branch that skips entire blocked blocks via \texttt{continue}, avoiding unnecessary dot products.

\subsection{Production Considerations}

\paragraph{Memory.} Metadata: 320\,KB for 65K blocks. Mask: 1\,MB for batch of $256 \times 4096$.
\paragraph{Thread safety.} DAI timing DFA uses per-tenant state with thread-safe locks.
\paragraph{Deployment.} TSAM is an extension module, not a fork---no changes to vLLM's core data structures, API, or scheduler.

% ==============================================================================
% SECTION 8 — EXPERIMENTAL EVALUATION
% ==============================================================================
\section{Experimental Evaluation}
\label{sec:experiments}

We evaluate the reference implementation along five axes: (1) isolation
correctness on the worst-case block-table-confusion scenario, (2) tenant
scalability, (3) prefix-sharing memory accounting, (4) timing-defense
behaviour under the analytic adversary model, and (5) reference-kernel
correctness.  The headline result---bitwise-exact noninterference up to
24 transformer layers---is in \S\ref{sec:e2e:empirical}; the experiments
in this section measure mechanisms that build on top of that.

\paragraph{Setup.}  Unless noted, single-attention-layer experiments use
$d_h = 128$, $H = 32$, $H_{\mathrm{kv}} = 8$ (GQA), page size $P = 16$, and
a block pool of 8{,}192 blocks.  Section~\ref{sec:e2e:empirical}
experiments use $d_{\mathrm{model}} = 256$, $H = 4$, depths
$L \in \{1, 2, 4, 8, 12, 24\}$.  Experiments are CPU-executable; full GPU
performance numbers (FlashAttention-2 fused kernel, vLLM continuous
batching, Llama-2 throughput) are deferred to a future evaluation
campaign with hardware-time and gated-model access (see
\S\ref{sec:limitations}).

\textbf{Baselines}: Vanilla attention (no isolation), Orthogonal
Projection~\citep{li2024orthogonal} ($d_{\mathrm{sub}} = 256$), Process
Isolation.

\subsection{Experiment 1: Isolation Correctness}

\textbf{RQ1: Does TSAM achieve zero leakage?}

\emph{Setup}: Worst-case scenario---attacker's block table includes victim's physical blocks. Victim KV content set to distinctive value. 100 attack queries.

\begin{table}[t]
\centering
\caption{Cross-tenant information leakage (bits per query). Lower is better.}
\label{tab:isolation}
\begin{tabular}{@{}lccc@{}}
\toprule
\textbf{Method} & \textbf{Bits/Query} & \textbf{Cross-Tenant Weight} & \textbf{Guarantee} \\
\midrule
Vanilla vLLM & 3.71 & $> 0$ & None \\
Orth.\ Projection & $\sim 10^{-6}$ & $\varepsilon > 0$ & Approximate \\
\textbf{TSAM (ours)} & \textbf{0.0} & \textbf{exactly 0.0} & \textbf{Exact ($I{=}0$)} \\
Process Isolation & 0.0 & 0.0 & Exact \\
\bottomrule
\end{tabular}
\end{table}

TSAM achieves 0.0 bits across all 100 queries. Every cross-tenant position is masked to $-\infty$, verified via direct mask inspection. TSAM matches process isolation at $<$1\% overhead vs.\ 5--10$\times$.

\subsection{Experiment 2: Scalability}

\textbf{RQ2: Does isolation hold at arbitrary scale?}

\begin{table}[t]
\centering
\caption{Scalability: leakage (bits) vs.\ tenant count. TSAM scales to 10,000+ tenants; Orthogonal Projection fails at $N > 16$.}
\label{tab:scalability}
\begin{tabular}{@{}rcccc@{}}
\toprule
\textbf{Tenants} & \textbf{TSAM Leakage} & \textbf{TSAM Status} & \textbf{Orth.\ Proj.\ Leakage} & \textbf{Orth.\ Proj.\ Status} \\
\midrule
10 & 0.0 & \cmark & 0.0 & \cmark \\
16 & 0.0 & \cmark & 0.0 & \cmark \\
17 & 0.0 & \cmark & --- & \textsc{fail} \\
100 & 0.0 & \cmark & --- & \textsc{fail} \\
1,000 & 0.0 & \cmark & --- & \textsc{fail} \\
10,000 & 0.0 & \cmark & --- & \textsc{fail} \\
\bottomrule
\end{tabular}
\end{table}

Orthogonal Projection fails at $N = 17$ ($17 \times 256 = 4352 > d_{\mathrm{model}} = 4096$). \textbf{TSAM is the first isolation mechanism achieving zero leakage at 10,000+ concurrent tenants.}

\subsection{Experiment 3: Prefix Sharing Efficiency}

\textbf{RQ3: Is prefix sharing preserved?}

100 tenants, 1,024-token shared prefix.

\begin{table}[t]
\centering
\caption{Prefix sharing: TSAM preserves single-copy sharing while maintaining isolation.}
\label{tab:prefix}
\begin{tabular}{@{}lcccc@{}}
\toprule
\textbf{Method} & \textbf{Prefix Copies} & \textbf{Memory} & \textbf{Throughput vs.\ Vanilla} & \textbf{Isolation} \\
\midrule
Vanilla & 1 & $1\times$ & 100\% (baseline) & \xmark \\
\textbf{TSAM} & \textbf{1} & $\mathbf{1\times}$ & \textbf{95--100\%} & \cmark \\
Orth.\ Proj. & $N$ (100) & $100\times$ & $\sim$60\% & \cmark$^*$ \\
Process Iso. & $N$ (100) & $100\times$ & $\sim$40\% & \cmark \\
\bottomrule
\multicolumn{5}{l}{\footnotesize $^*$Limited to $N \leq 16$; extrapolated.}
\end{tabular}
\end{table}

TSAM stores one prefix copy (32\,MB) vs.\ Orthogonal Projection's 100 copies (3.2\,GB)---a $100\times$ memory reduction. TSAM is the \textbf{only} method achieving isolation with single-copy prefix sharing.

\subsection{Experiment 4: Timing Side-Channel Defense}

\textbf{RQ4: Does DAI neutralize timing attacks?}

1,000 probes; cache hit base $= 30$\,ms ($\sigma = 5$), miss base $= 80$\,ms ($\sigma = 10$).

\begin{table}[t]
\centering
\caption{Timing defense: DAI reduces attacker accuracy to near-random guessing.}
\label{tab:timing}
\begin{tabular}{@{}lccc@{}}
\toprule
\textbf{Defense} & \textbf{Accuracy} & \textbf{Bits/Query} & \textbf{Reduction} \\
\midrule
No defense & $\sim$97\% & $\sim$0.95 & --- \\
Noise only & $\sim$75\% & $\sim$0.40 & 58\% \\
\textbf{TSAM + DAI} & $\mathbf{\sim}$\textbf{52\%} & $\mathbf{< 0.05}$ & $\mathbf{\geq 95\%}$ \\
\bottomrule
\end{tabular}
\end{table}

TSAM+DAI reduces accuracy to 52\%---statistically indistinguishable from random guessing (50\%).

\subsection{Experiment 5: Throughput Overhead (Production Gap)}
\label{sec:experiments:throughput}

\textbf{RQ5: What is the performance cost?}

The TSAM mask is a metadata lookup applied before softmax: per-position
cost is $\Theta(1)$, dominated by a uint8 comparison and an int32
comparison, both warp-uniform.  Analytically, the per-block instruction
overhead is $\approx$3 instructions, amortized across the full
$P \times d_h$ multiply-accumulate of the dot-product
($\approx 0.015\%$ of dot-product compute, before any kernel-fusion
optimisation).

End-to-end production throughput numbers---FlashAttention-2 fused with the
TSAM mask, deployed in a real vLLM fork running Llama-2-7B/13B with
continuous batching on A100/H100---are not yet available in this version
and would not be defensible to claim from a Triton-prototype measurement.
We catalogue this as the single largest empirical gap between the present
work and a deployment-ready submission.  The reproducibility runbook in
\texttt{BENCHMARKS.md} specifies the exact command sequence that will
generate Table~\ref{tab:throughput-todo} once the GPU benchmark campaign
is complete.

\begin{table}[t]
\centering
\caption{Throughput overhead vs.\ vanilla attention.  Numbers to be
measured on production GPU under FlashAttention-2 + vLLM; until the
benchmark campaign completes, this table is a runbook
specification, not a result.}
\label{tab:throughput-todo}
\begin{tabular}{@{}lrrr@{}}
\toprule
\textbf{Config ($B \times S$, model)} & \textbf{Vanilla (tok/s)} & \textbf{TSAM (tok/s)} & \textbf{Overhead} \\
\midrule
$1 \times 512$, Llama-2-7B & [TODO: GB10] & [TODO: GB10] & [TODO] \\
$8 \times 2{,}048$, Llama-2-7B & [TODO: A100] & [TODO: A100] & [TODO] \\
$32 \times 8{,}192$, Llama-2-7B & [TODO: A100] & [TODO: A100] & [TODO] \\
$128 \times 32{,}768$, Llama-2-13B & [TODO: H100] & [TODO: H100] & [TODO] \\
\bottomrule
\end{tabular}
\end{table}

\subsection{Experiment 6: Kernel Correctness}

Systematic verification across all configurations:
\begin{enumerate}[label=(\alph*),leftmargin=*,itemsep=2pt]
    \item \textbf{Own-tenant PRIVATE}: all positions unmasked, output non-zero \cmark
    \item \textbf{Cross-tenant PRIVATE}: all positions masked, weight \emph{exactly} 0.0 \cmark
    \item \textbf{SHARED pages}: accessible by all tenants \cmark
    \item \textbf{Edge cases}: batch 1/32/512, seq len 1/512/32768, page boundaries \cmark
\end{enumerate}
All tests pass. No numerical precision issues arise because TSAM operates on integer metadata comparisons, not floating-point projections.


% ==============================================================================
% SECTION 9 — DISCUSSION
% ==============================================================================
\section{Discussion}
\label{sec:discussion}

\subsection{Why Hard Masking Beats Soft Isolation}

\paragraph{Exact versus approximate isolation.}
Hard masking sets cross-tenant attention weights to \emph{structural zero} via $\exp(-\infty) = 0$. Soft isolation (orthogonal projection) achieves only numerical near-zero ($\sim 10^{-6}$), leaving a residual channel that compounds across layers~\citep{li2024orthogonal}. Recent extraction attacks~\citep{carlini2021extracting, nasr2023scalable} demonstrate that even small information channels can be exploited at scale.

\paragraph{Decoupled from dimensionality.}
TSAM's mask is a metadata lookup---$O(n)$ independent of $d_{\mathrm{model}}$. Orthogonal projection requires $N \times d_{\mathrm{sub}} \leq d_{\mathrm{model}}$, limiting tenants to $\lfloor d_{\mathrm{model}}/d_{\mathrm{sub}} \rfloor = 16$. TSAM imposes no such ceiling.

\paragraph{Preservation of prefix sharing.}
Hard masking operates on attention logits \emph{after} dot products, leaving shared KV-cache content untransformed. All tenants attend to identical shared pages. Projection transforms shared content per-tenant, destroying sharing.

\paragraph{Key insight.} \emph{Isolation is a data-flow property best enforced at the attention level, not the representation level.} This aligns with the noninterference framework of \citet{goguen1982security}: control which inputs influence outputs, rather than transforming inputs.

\subsection{Limitations}
\label{sec:limitations}

We are explicit about what TSAM does \emph{not} guarantee.  Each item below
is a real channel that a sufficiently motivated adversary can exploit;
deployment teams must reason about each one independently.

\paragraph{Channels out of scope.}
\begin{itemize}[leftmargin=1.5em,itemsep=2pt]
\item \textbf{GPU memory side channels.}  L2-cache timing, DRAM
  row-buffer conflicts, and CUDA scheduling patterns can leak page
  occupancy and access timing.  These require hardware/OS mitigations
  (e.g., spatial partitioning, MIG, microcode updates against
  Spectre~\citep{kocher2019spectre}, page-coloring), not software
  masking.  TSAM does not address them.
\item \textbf{Allocation-pattern metadata.}  The set of physical blocks a
  tenant holds, the timing of allocations and frees, and the sequence
  lengths processed all leak information about tenant activity and
  prompt characteristics.  TSAM does not obfuscate allocation metadata.
\item \textbf{Malicious operator.}  TSAM trusts the serving infrastructure
  to assign correct page types, apply the mask before softmax, and
  preserve metadata across continuous batching.  A malicious or
  compromised operator can trivially bypass it (e.g., by mislabelling
  a private page as shared).
\item \textbf{Non-attention cross-position channels.}  Architectures that
  perform cross-position operations outside attention---for example,
  sequence-level batch normalisation, 1D convolutions over positions, or
  cross-attention to external memory---violate condition C3 and are not
  covered by Theorem~\ref{thm:e2e}.  Standard decoder-only LMs (GPT,
  LLaMA, Qwen, Gemma) satisfy C3 by construction.
\item \textbf{Model-mediated semantic channels.}  TSAM controls
  information flow through attention weights.  It does not prevent an
  adversary from inferring information through the model's generated
  text via prompt injection, jailbreak prompts, or shared-prompt
  contamination.  These require application-layer defences.
\item \textbf{Training-data extraction.}  TSAM is an inference-time
  mechanism.  Memorisation in model weights~\citep{carlini2021extracting,
  nasr2023scalable, lukas2023analyzing} is orthogonal and unaddressed.
\item \textbf{Per-tenant LoRA adapters in the masked path.}  Per-tenant
  adapters~\citep{hu2022lora, sheng2023slora} placed inside the masked
  attention computation violate condition C4 and break the proof.
  Adapter-aware variants of TSAM are future work.
\end{itemize}

\paragraph{Empirical gaps in this paper.}
The strongest result we report is bitwise exact noninterference up to
24-layer transformers with random weights
(\S\ref{sec:e2e:empirical}).  Several measurements that a
deployment-ready submission would include are deferred:
\begin{itemize}[leftmargin=1.5em,itemsep=2pt]
\item \textbf{Real LLM serving throughput} (Llama-2-7B/13B on A100/H100
  under continuous batching) is catalogued as
  Table~\ref{tab:throughput-todo} but not yet measured.
\item \textbf{Real adversarial co-tenant harness} (block-table confusion,
  shared-prefix poisoning, stale-block reuse) under TSAM-vLLM is not
  yet runnable end-to-end; production timing-attack ROC/AUC requires
  this harness.
\item \textbf{Compiled CUDA kernel patch.}  The patch returned by
  \texttt{generate\_vllm\_kernel\_patch} is a sketch.  It uses
  \texttt{-INFINITY} for IEEE~754 alignment with the Python reference
  and the masking branch is in scope of the \texttt{qk} declaration, but
  it has not yet been compiled or applied to a pinned vLLM commit.
\end{itemize}

\paragraph{Detection window of the DAI defence.}
The DAI timing DFA is detection-triggered, not constant-time.  Pre-detection,
the adversary operates against an undefended channel; the maximum
pre-detection leakage is bounded by $K + W \leq 70$ probes (\S\ref{sec:dai}).
A constant-time padding variant would close this window at the cost of
unconditional latency overhead and is left as future work.

\paragraph{Production hardening.}  Production deployment additionally
requires: a clean-compute provenance check that verifies condition C2
before admitting any page as SHARED (the
\texttt{tsam.core.compute\_provenance} module sketches the API); a typed
page-lifecycle module that zeroizes freed KV slots before reuse and
enforces copy-on-write when a private continuation derives from a shared
prefix (\texttt{tsam.core.page\_lifecycle}); and an admission policy on
shared-prefix content hashes so that an attacker cannot launder private
data through the SHARED label.

\subsection{Broader Impact}

TSAM enables secure shared LLM infrastructure with 10--100$\times$ cost reduction vs.\ process isolation. The formal guarantees---grounded in noninterference~\citep{goguen1982security} and information-theoretic zero leakage~\citep{shannon1948mathematical}---provide a basis for compliance with HIPAA, SOC 2, and FedRAMP. All code is released open-source for reproducibility.

% ==============================================================================
% SECTION 10 — CONCLUSION
% ==============================================================================
\section{Conclusion}
\label{sec:conclusion}

We identified a fundamental gap in multi-tenant LLM serving: \emph{no existing approach achieves zero leakage, unlimited tenant scale, and prefix sharing simultaneously}. Process isolation guarantees safety but kills efficiency; orthogonal projection scales only to 16 tenants and destroys prefix sharing; ad hoc masks lack formal guarantees.

TSAM closes this gap through \emph{typed KV-cache pages} and \emph{hard attention masking}. Each cache page carries a lightweight type annotation (PRIVATE or SHARED), and a single conditional in the attention kernel enforces:
\begin{equation}
    I\!\bigl(\mathrm{output}(q_i);\;\mathrm{KV}_{\mathrm{private}}(t_j)\bigr) = 0 \quad \forall\, i \neq j.
\end{equation}

Our evaluation demonstrates:
\begin{itemize}[leftmargin=*,itemsep=2pt]
    \item \textbf{Bitwise-exact end-to-end noninterference} on a 24-layer
      transformer: the attacker's output logits are bitwise-identical
      regardless of the victim's private tokens, across 50 trials per
      depth $L \in \{1,2,4,8,12,24\}$ (\S\ref{sec:e2e:empirical}).
    \item \textbf{Linear probing collapse}: a ridge classifier extracts
      the victim's private prompt class with 100\% accuracy under
      vanilla attention but only random chance ($11.7\%$ vs.\ $12.5\%$
      baseline) under TSAM, at every depth tested.
    \item \textbf{Memory-bound tenant capacity}: the mask is a metadata
      lookup with no algebraic dependence on $d_{\mathrm{model}}$ or
      tenant count; orthogonal projection fails at
      $\lfloor d_{\mathrm{model}}/d_{\mathrm{sub}} \rfloor$ tenants by
      construction (\S\ref{sec:formal:orthogonal}).
    \item \textbf{Preserved prefix sharing}: SHARED pages remain
      single-copy and visible to all tenants (\S\ref{sec:design:pages}).
    \item \textbf{Per-page metadata}: 5 bytes per block, $\approx$0.06\%
      of KV-cache volume.
    \item \textbf{Detection-triggered timing mitigation}: DAI bounds
      pre-detection leakage to $\leq K + W = 70$ probes and reduces
      post-detection per-probe channel capacity by $>99\%$ under the
      analytic adversary model (\S\ref{sec:dai}).
\end{itemize}

We are explicit about what we do not claim: production GPU
throughput numbers (Table~\ref{tab:throughput-todo} is a runbook
specification, not a result), end-to-end TSAM-vLLM integration on
Llama-2 weights, and protection against channels outside attention
(\S\ref{sec:limitations}).

The key insight is that \emph{isolation is a data-flow property best
enforced at the attention level, not the representation level}.  By
intervening where information flows between tokens---rather than
transforming the representations themselves---TSAM achieves structural
zero-leakage in the attention channel, decouples tenant scalability
from model dimensionality, and preserves prefix sharing.

\paragraph{Future work.}  The deployment-ready submission this paper
points toward requires: (i) a real vLLM fork with the CUDA patch
compiled and applied, (ii) Llama-2-7B/13B end-to-end throughput on
A100/H100 with continuous batching, (iii) a real adversarial co-tenant
harness exercising block-table confusion and shared-prefix poisoning,
(iv) a constant-time variant of DAI to close the detection window, and
(v) extending TSAM to per-tenant LoRA adapters~\citep{hu2022lora,
sheng2023slora}, distributed multi-GPU serving, and formal verification
in Coq or Lean.

% ==============================================================================

\bibliographystyle{plainnat}
\bibliography{references}

\end{document}
